\begin{explanationbox}
	\textbf{\textcolor{CExplanation}{一句话直觉}}

	二次函数图像是抛物线，顶点公式直接给出最高或最低点的位置，对称轴则把抛物线分成左右对称的两部分。

	\medskip
	\textbf{\textcolor{CExplanation}{核心拆解}}
	\begin{description}[style=nextline, leftmargin=2.8em, labelsep=0.5em, itemsep=0.45em]
		\item[\textbf{要点一（配方变形）：}] 将一般式 $y=ax^2+bx+c$ 通过配方法写成 $a(x-h)^2+k$ 的形式，其中 $h=-\dfrac{b}{2a}$，$k=\dfrac{4ac-b^2}{4a}$。
		\item[\textbf{要点二（顶点来源）：}] 平方项 $a(x-h)^2$ 的最小值为 $0$（$a>0$）或最大值为 $0$（$a<0$），顶点在 $x=h$ 处取得极值。
		\item[\textbf{要点三（对称轴关系）：}] 顶点横坐标 $h$ 恰好是抛物线对称轴的位置，纵坐标 $k$ 是抛物线最高或最低点的值。
	\end{description}

	\medskip
	\textbf{\textcolor{CExplanation}{几何本质（对称与最值）}}

	抛物线的顶点是曲线上的唯一极值点，对称轴经过顶点且平分抛物线。

	\medskip
	\textbf{\textcolor{CExplanation}{代数意义（配方法核心）}}

	配方将一般式转化为顶点式，直接揭示了函数的极值点、对称轴和函数值的范围。

	\medskip
	\textbf{\textcolor{CExplanation}{考点价值（求极值与图像）}}
	\begin{itemize}[leftmargin=*, itemsep=0.5em, topsep=0.4em]
		\item \textbf{考法一（求顶点坐标）：} 给出一般式，直接套用公式得顶点坐标，无需配方。
		\item \textbf{考法二（求最值范围）：} 由顶点纵坐标及开口方向，确定函数的最大值或最小值。
	\end{itemize}

	\medskip
	\textbf{\textcolor{CExplanation}{顿悟点}}

	横坐标公式是$x=-\dfrac{b}{2a}$，纵坐标公式是 $\dfrac{4ac-b^2}{4a}$，两者都是由系数唯一决定，无需死记硬背配方过程。

	\medskip
	\textbf{\textcolor{CExplanation}{使用场景}}
	\begin{itemize}[leftmargin=*, itemsep=0.5em, topsep=0.4em]
		\item \textbf{场景一（画抛物线草图）：} 先由顶点公式确定顶点，再由对称轴和开口画抛物线。
		\item \textbf{场景二（最值应用题）：} 利润、面积等二次函数模型，直接利用顶点纵坐标求最值。
	\end{itemize}
\end{explanationbox}
